% --- Archon physics formalization source begin ---
% archon:physics
% archon:covers PhyXMiniProblems/problem_phyx_mini_0823.lean
% archon:source-report reports/phyx_mini/problem_phyx_mini_0823.source.json

\chapter{Physics problem phyx\_mini\_0823}
\label{ch:PhyXMiniProblems_problem_phyx_mini_0823}

\paragraph{Problem source.}
\#\# Physical scenario

To loosen a pipe fitting, a plumber slips a piece of scrap pipe (a ``cheater'') over his wrench handle. He stands on the end of the cheater, applying his 900 N weight at a point 0.80 m from the center of the fitting as shown in figure. The wrench handle and cheater make an angle of \textbackslash{}( 19\textasciicircum{}\textbackslash{}circ \textbackslash{}) with the horizontal.

\#\# Question

Find the magnitude and direction of the torque he applies about the center of the fitting.

\#\# Answer choices

- A: A: \textbackslash{}( 700 \textbackslash{}text\{ N\} \textbackslash{}cdot \textbackslash{}text\{ m\} \textbackslash{})

- B: B: \textbackslash{}( 680 \textbackslash{}text\{ N\} \textbackslash{}cdot \textbackslash{}text\{ m\} \textbackslash{})

- C: C: \textbackslash{}( 340 \textbackslash{}text\{ N\} \textbackslash{}cdot \textbackslash{}text\{ m\} \textbackslash{})

- D: D: \textbackslash{}( 350 \textbackslash{}text\{ N\} \textbackslash{}cdot \textbackslash{}text\{ m\} \textbackslash{})

\#\# Figure caption (auxiliary; use the image as primary evidence)

The image depicts a mechanical scenario involving a person applying force to a wrench. Here are the details:

- A person is using their foot to apply a downward force on the handle of a pipe wrench.
- The magnitude of the force applied is labeled as \textbackslash{}( F = 900 \textbackslash{}, \textbackslash{}text\{N\} \textbackslash{}).
- The direction of the force is shown with a red arrow pointing straight down.
- The length of the wrench handle from the point of force application to the pivot point is marked as \textbackslash{}( 0.80 \textbackslash{}, \textbackslash{}text\{m\} \textbackslash{}).
- The wrench is positioned at an angle, labeled as \textbackslash{}( 19\textasciicircum{}\textbackslash{}circ \textbackslash{}) above the horizontal.
- The pivot point of the wrench is resting on a grey pipe, indicating where the wrench is turning.
- The wrench is orange, and the dashed line represents the horizontal for reference to the angle measurement.

\#\# Dataset metadata

Rotational Motion; Physical Model Grounding Reasoning, Spatial Relation Reasoning

\paragraph{Recorded answer/context.}
B: B: \textbackslash{}( 680 \textbackslash{}text\{ N\} \textbackslash{}cdot \textbackslash{}text\{ m\} \textbackslash{})

\paragraph{Figure/image path.}
/root/proposal\_for\_physic/science-mango/phyx\_mini\_run/phyx\_data/test\_image/823.png

\paragraph{Formalization target.}
create a compiling Lean file with sorry bodies at `PhyXMiniProblems/problem\_phyx\_mini\_0823.lean`. The Lean declarations must preserve the physical quantities, dimensions or dimensional roles, figure labels, governing-law hypotheses, and final relation expressed by this problem.
Use Mathlib/Physlib names found through LeanExplore where available. If a physics API is missing, introduce faithful local abstractions rather than scalar placeholder aliases.

\begin{theorem}[Physics formalization target]
\label{thm:physics:phyx_mini_0823:target}
\lean{PhyXMiniProblems.ProblemPhyXMini0823.problem_phyx_mini_0823}
\uses{def:physics:phyx-mini-0823:phyxminiproblems-problemphyxmini0823-torqueinnewtonmeters, def:physics:phyx-mini-0823:phyxminiproblems-problemphyxmini0823-torquemagnitudeinnewtonmeters, def:physics:phyx-mini-0823:phyxminiproblems-problemphyxmini0823-zhat, def:physics:phyx-mini-0823:phyxminiproblems-problemphyxmini0823-pipefittingtorquesetup, def:physics:phyx-mini-0823:phyxminiproblems-problemphyxmini0823-matchesproblemstatement, def:physics:phyx-mini-0823:phyxminiproblems-problemphyxmini0823-matchesprimaryfigure, def:physics:phyx-mini-0823:phyxminiproblems-problemphyxmini0823-matcheshandlegeometry, def:physics:phyx-mini-0823:phyxminiproblems-problemphyxmini0823-matchesdownwardforcegeometry, def:physics:phyx-mini-0823:phyxminiproblems-problemphyxmini0823-satisfiestorquelaw, def:physics:phyx-mini-0823:phyxminiproblems-problemphyxmini0823-recordedanswerchoice, def:physics:phyx-mini-0823:phyxminiproblems-problemphyxmini0823-agreesattennewtonmeterprecision, def:physics:phyx-mini-0823:phyxminiproblems-problemphyxmini0823-torquesense-value}
The assigned autoformalize agent should translate this physics problem into Lean declarations in the covered file, with theorem and lemma proof bodies written as `by sorry`.
\end{theorem}
\begin{proof}
This is an autoformalization task, not a proof task. Produce statements that can later be proved by the physics prover without weakening the physical model.
\end{proof}

% --- Archon declaration topology begin ---
\section{Declaration topology}
\begin{definition}[Spatial Vector]
  \label{def:physics:phyx-mini-0823:phyxminiproblems-problemphyxmini0823-spatialvector}
  \lean{PhyXMiniProblems.ProblemPhyXMini0823.SpatialVector}
  Three-dimensional Euclidean vectors used for coherent-unit readouts
\end{definition}
\begin{proof}
  This is the declared model definition of Spatial Vector.
\end{proof}
\begin{definition}[force Dimension]
  \label{def:physics:phyx-mini-0823:phyxminiproblems-problemphyxmini0823-forcedimension}
  \lean{PhyXMiniProblems.ProblemPhyXMini0823.forceDimension}
  Force has physical dimension M L T⁻²
\end{definition}
\begin{proof}
  This is the declared model definition of force Dimension.
\end{proof}
\begin{definition}[torque Dimension]
  \label{def:physics:phyx-mini-0823:phyxminiproblems-problemphyxmini0823-torquedimension}
  \lean{PhyXMiniProblems.ProblemPhyXMini0823.torqueDimension}
  \uses{def:physics:phyx-mini-0823:phyxminiproblems-problemphyxmini0823-forcedimension}
  Torque has physical dimension M L² T⁻²
\end{definition}
\begin{proof}
  This is the declared model definition of torque Dimension.
\end{proof}
\begin{definition}[Length Quantity]
  \label{def:physics:phyx-mini-0823:phyxminiproblems-problemphyxmini0823-lengthquantity}
  \lean{PhyXMiniProblems.ProblemPhyXMini0823.LengthQuantity}
  A unit-independent nonnegative physical length
\end{definition}
\begin{proof}
  This is the declared model definition of Length Quantity.
\end{proof}
\begin{definition}[Position Vector Quantity]
  \label{def:physics:phyx-mini-0823:phyxminiproblems-problemphyxmini0823-positionvectorquantity}
  \lean{PhyXMiniProblems.ProblemPhyXMini0823.PositionVectorQuantity}
  \uses{def:physics:phyx-mini-0823:phyxminiproblems-problemphyxmini0823-spatialvector}
  A unit-independent spatial position with length dimension
\end{definition}
\begin{proof}
  This is the declared model definition of Position Vector Quantity.
\end{proof}
\begin{definition}[Force Vector Quantity]
  \label{def:physics:phyx-mini-0823:phyxminiproblems-problemphyxmini0823-forcevectorquantity}
  \lean{PhyXMiniProblems.ProblemPhyXMini0823.ForceVectorQuantity}
  \uses{def:physics:phyx-mini-0823:phyxminiproblems-problemphyxmini0823-spatialvector, def:physics:phyx-mini-0823:phyxminiproblems-problemphyxmini0823-forcedimension}
  A unit-independent spatial force vector
\end{definition}
\begin{proof}
  This is the declared model definition of Force Vector Quantity.
\end{proof}
\begin{definition}[Torque Vector Quantity]
  \label{def:physics:phyx-mini-0823:phyxminiproblems-problemphyxmini0823-torquevectorquantity}
  \lean{PhyXMiniProblems.ProblemPhyXMini0823.TorqueVectorQuantity}
  \uses{def:physics:phyx-mini-0823:phyxminiproblems-problemphyxmini0823-spatialvector, def:physics:phyx-mini-0823:phyxminiproblems-problemphyxmini0823-torquedimension}
  A unit-independent spatial torque vector about the fitting center
\end{definition}
\begin{proof}
  This is the declared model definition of Torque Vector Quantity.
\end{proof}
\begin{definition}[nonnegative Readout]
  \label{def:physics:phyx-mini-0823:phyxminiproblems-problemphyxmini0823-nonnegativereadout}
  \lean{PhyXMiniProblems.ProblemPhyXMini0823.nonnegativeReadout}
  Read a nonnegative physical scalar in a coherent unit system
\end{definition}
\begin{proof}
  This is the declared model definition of nonnegative Readout.
\end{proof}
\begin{definition}[vector Readout]
  \label{def:physics:phyx-mini-0823:phyxminiproblems-problemphyxmini0823-vectorreadout}
  \lean{PhyXMiniProblems.ProblemPhyXMini0823.vectorReadout}
  \uses{def:physics:phyx-mini-0823:phyxminiproblems-problemphyxmini0823-spatialvector}
  Read a physical vector in a coherent unit system
\end{definition}
\begin{proof}
  This is the declared model definition of vector Readout.
\end{proof}
\begin{definition}[length In Meters]
  \label{def:physics:phyx-mini-0823:phyxminiproblems-problemphyxmini0823-lengthinmeters}
  \lean{PhyXMiniProblems.ProblemPhyXMini0823.lengthInMeters}
  \uses{def:physics:phyx-mini-0823:phyxminiproblems-problemphyxmini0823-lengthquantity, def:physics:phyx-mini-0823:phyxminiproblems-problemphyxmini0823-nonnegativereadout}
  Metre readout of a physical length
\end{definition}
\begin{proof}
  This is the declared model definition of length In Meters.
\end{proof}
\begin{definition}[position In Meters]
  \label{def:physics:phyx-mini-0823:phyxminiproblems-problemphyxmini0823-positioninmeters}
  \lean{PhyXMiniProblems.ProblemPhyXMini0823.positionInMeters}
  \uses{def:physics:phyx-mini-0823:phyxminiproblems-problemphyxmini0823-spatialvector, def:physics:phyx-mini-0823:phyxminiproblems-problemphyxmini0823-positionvectorquantity, def:physics:phyx-mini-0823:phyxminiproblems-problemphyxmini0823-vectorreadout}
  Cartesian metre readout of a position
\end{definition}
\begin{proof}
  This is the declared model definition of position In Meters.
\end{proof}
\begin{definition}[force In Newtons]
  \label{def:physics:phyx-mini-0823:phyxminiproblems-problemphyxmini0823-forceinnewtons}
  \lean{PhyXMiniProblems.ProblemPhyXMini0823.forceInNewtons}
  \uses{def:physics:phyx-mini-0823:phyxminiproblems-problemphyxmini0823-spatialvector, def:physics:phyx-mini-0823:phyxminiproblems-problemphyxmini0823-forcevectorquantity, def:physics:phyx-mini-0823:phyxminiproblems-problemphyxmini0823-vectorreadout}
  Cartesian newton readout of a force
\end{definition}
\begin{proof}
  This is the declared model definition of force In Newtons.
\end{proof}
\begin{definition}[force Magnitude In Newtons]
  \label{def:physics:phyx-mini-0823:phyxminiproblems-problemphyxmini0823-forcemagnitudeinnewtons}
  \lean{PhyXMiniProblems.ProblemPhyXMini0823.forceMagnitudeInNewtons}
  \uses{def:physics:phyx-mini-0823:phyxminiproblems-problemphyxmini0823-forcevectorquantity, def:physics:phyx-mini-0823:phyxminiproblems-problemphyxmini0823-forceinnewtons}
  Euclidean magnitude, in newtons, of a physical force
\end{definition}
\begin{proof}
  This is the declared model definition of force Magnitude In Newtons.
\end{proof}
\begin{definition}[torque In Newton Meters]
  \label{def:physics:phyx-mini-0823:phyxminiproblems-problemphyxmini0823-torqueinnewtonmeters}
  \lean{PhyXMiniProblems.ProblemPhyXMini0823.torqueInNewtonMeters}
  \uses{def:physics:phyx-mini-0823:phyxminiproblems-problemphyxmini0823-spatialvector, def:physics:phyx-mini-0823:phyxminiproblems-problemphyxmini0823-torquevectorquantity, def:physics:phyx-mini-0823:phyxminiproblems-problemphyxmini0823-vectorreadout}
  Cartesian newton-metre readout of a torque
\end{definition}
\begin{proof}
  This is the declared model definition of torque In Newton Meters.
\end{proof}
\begin{definition}[torque Magnitude In Newton Meters]
  \label{def:physics:phyx-mini-0823:phyxminiproblems-problemphyxmini0823-torquemagnitudeinnewtonmeters}
  \lean{PhyXMiniProblems.ProblemPhyXMini0823.torqueMagnitudeInNewtonMeters}
  \uses{def:physics:phyx-mini-0823:phyxminiproblems-problemphyxmini0823-torquevectorquantity, def:physics:phyx-mini-0823:phyxminiproblems-problemphyxmini0823-torqueinnewtonmeters}
  Euclidean magnitude, in newton-metres, of a physical torque
\end{definition}
\begin{proof}
  This is the declared model definition of torque Magnitude In Newton Meters.
\end{proof}
\begin{definition}[spatial Cross]
  \label{def:physics:phyx-mini-0823:phyxminiproblems-problemphyxmini0823-spatialcross}
  \lean{PhyXMiniProblems.ProblemPhyXMini0823.spatialCross}
  \uses{def:physics:phyx-mini-0823:phyxminiproblems-problemphyxmini0823-spatialvector}
  The ordinary three-dimensional cross product transported to Euclidean space
\end{definition}
\begin{proof}
  This is the declared model definition of spatial Cross.
\end{proof}
\begin{definition}[x Hat]
  \label{def:physics:phyx-mini-0823:phyxminiproblems-problemphyxmini0823-xhat}
  \lean{PhyXMiniProblems.ProblemPhyXMini0823.xHat}
  \uses{def:physics:phyx-mini-0823:phyxminiproblems-problemphyxmini0823-spatialvector}
  Dimensionless Cartesian unit vector pointing rightward in the image
\end{definition}
\begin{proof}
  This is the declared model definition of x Hat.
\end{proof}
\begin{definition}[y Hat]
  \label{def:physics:phyx-mini-0823:phyxminiproblems-problemphyxmini0823-yhat}
  \lean{PhyXMiniProblems.ProblemPhyXMini0823.yHat}
  \uses{def:physics:phyx-mini-0823:phyxminiproblems-problemphyxmini0823-spatialvector}
  Dimensionless Cartesian unit vector pointing upward in the image
\end{definition}
\begin{proof}
  This is the declared model definition of y Hat.
\end{proof}
\begin{definition}[z Hat]
  \label{def:physics:phyx-mini-0823:phyxminiproblems-problemphyxmini0823-zhat}
  \lean{PhyXMiniProblems.ProblemPhyXMini0823.zHat}
  \uses{def:physics:phyx-mini-0823:phyxminiproblems-problemphyxmini0823-spatialvector}
  Dimensionless Cartesian unit vector pointing out of the page
\end{definition}
\begin{proof}
  This is the declared model definition of z Hat.
\end{proof}
\begin{definition}[Wrench Point]
  \label{def:physics:phyx-mini-0823:phyxminiproblems-problemphyxmini0823-wrenchpoint}
  \lean{PhyXMiniProblems.ProblemPhyXMini0823.WrenchPoint}
  The two distinguished points needed for the lever arm about the fitting
\end{definition}
\begin{proof}
  This is the declared model definition of Wrench Point.
\end{proof}
\begin{definition}[Figure Component]
  \label{def:physics:phyx-mini-0823:phyxminiproblems-problemphyxmini0823-figurecomponent}
  \lean{PhyXMiniProblems.ProblemPhyXMini0823.FigureComponent}
  Individually visible components of the supplied bitmap
\end{definition}
\begin{proof}
  This is the declared model definition of Figure Component.
\end{proof}
\begin{definition}[Figure Label]
  \label{def:physics:phyx-mini-0823:phyxminiproblems-problemphyxmini0823-figurelabel}
  \lean{PhyXMiniProblems.ProblemPhyXMini0823.FigureLabel}
  Literal calibrated labels printed in the supplied bitmap
\end{definition}
\begin{proof}
  This is the declared model definition of Figure Label.
\end{proof}
\begin{definition}[Vertical Direction]
  \label{def:physics:phyx-mini-0823:phyxminiproblems-problemphyxmini0823-verticaldirection}
  \lean{PhyXMiniProblems.ProblemPhyXMini0823.VerticalDirection}
  Qualitative direction of the applied-force arrow
\end{definition}
\begin{proof}
  This is the declared model definition of Vertical Direction.
\end{proof}
\begin{definition}[Handle Extension]
  \label{def:physics:phyx-mini-0823:phyxminiproblems-problemphyxmini0823-handleextension}
  \lean{PhyXMiniProblems.ProblemPhyXMini0823.HandleExtension}
  The method used to extend the wrench handle
\end{definition}
\begin{proof}
  This is the declared model definition of Handle Extension.
\end{proof}
\begin{definition}[Torque Sense]
  \label{def:physics:phyx-mini-0823:phyxminiproblems-problemphyxmini0823-torquesense}
  \lean{PhyXMiniProblems.ProblemPhyXMini0823.TorqueSense}
  Rotation sense in the side view, with positive z pointing out of the page
\end{definition}
\begin{proof}
  This is the declared model definition of Torque Sense.
\end{proof}
\begin{definition}[Supplied Wrench Figure]
  \label{def:physics:phyx-mini-0823:phyxminiproblems-problemphyxmini0823-suppliedwrenchfigure}
  \lean{PhyXMiniProblems.ProblemPhyXMini0823.SuppliedWrenchFigure}
  \uses{def:physics:phyx-mini-0823:phyxminiproblems-problemphyxmini0823-wrenchpoint, def:physics:phyx-mini-0823:phyxminiproblems-problemphyxmini0823-figurecomponent, def:physics:phyx-mini-0823:phyxminiproblems-problemphyxmini0823-figurelabel, def:physics:phyx-mini-0823:phyxminiproblems-problemphyxmini0823-verticaldirection}
  Literal labels and incidences transcribed from image 823.png
\end{definition}
\begin{proof}
  This is the declared model definition of Supplied Wrench Figure.
\end{proof}
\begin{definition}[Pipe Fitting Torque Setup]
  \label{def:physics:phyx-mini-0823:phyxminiproblems-problemphyxmini0823-pipefittingtorquesetup}
  \lean{PhyXMiniProblems.ProblemPhyXMini0823.PipeFittingTorqueSetup}
  \uses{def:physics:phyx-mini-0823:phyxminiproblems-problemphyxmini0823-lengthquantity, def:physics:phyx-mini-0823:phyxminiproblems-problemphyxmini0823-positionvectorquantity, def:physics:phyx-mini-0823:phyxminiproblems-problemphyxmini0823-forcevectorquantity, def:physics:phyx-mini-0823:phyxminiproblems-problemphyxmini0823-torquevectorquantity, def:physics:phyx-mini-0823:phyxminiproblems-problemphyxmini0823-wrenchpoint, def:physics:phyx-mini-0823:phyxminiproblems-problemphyxmini0823-handleextension, def:physics:phyx-mini-0823:phyxminiproblems-problemphyxmini0823-suppliedwrenchfigure}
  Independent physical quantities in the wrench experiment. In particular, torqueAboutFitting is not defined from a displayed answer; the governing law below constrains it
\end{definition}
\begin{proof}
  This is the declared model definition of Pipe Fitting Torque Setup.
\end{proof}
\begin{definition}[Matches Problem Statement]
  \label{def:physics:phyx-mini-0823:phyxminiproblems-problemphyxmini0823-matchesproblemstatement}
  \lean{PhyXMiniProblems.ProblemPhyXMini0823.MatchesProblemStatement}
  \uses{def:physics:phyx-mini-0823:phyxminiproblems-problemphyxmini0823-lengthinmeters, def:physics:phyx-mini-0823:phyxminiproblems-problemphyxmini0823-forcemagnitudeinnewtons, def:physics:phyx-mini-0823:phyxminiproblems-problemphyxmini0823-pipefittingtorquesetup}
  Numerical and categorical data explicitly stated in the problem prose
\end{definition}
\begin{proof}
  This is the declared model definition of Matches Problem Statement.
\end{proof}
\begin{definition}[Matches Primary Figure]
  \label{def:physics:phyx-mini-0823:phyxminiproblems-problemphyxmini0823-matchesprimaryfigure}
  \lean{PhyXMiniProblems.ProblemPhyXMini0823.MatchesPrimaryFigure}
  \uses{def:physics:phyx-mini-0823:phyxminiproblems-problemphyxmini0823-lengthinmeters, def:physics:phyx-mini-0823:phyxminiproblems-problemphyxmini0823-forcemagnitudeinnewtons, def:physics:phyx-mini-0823:phyxminiproblems-problemphyxmini0823-pipefittingtorquesetup}
  Primary-image evidence: visible objects, all three numerical labels, the downward force arrow at the left endpoint, and the up-left placement of that endpoint relative to the fitting. No torque value or answer choice occurs here
\end{definition}
\begin{proof}
  This is the declared model definition of Matches Primary Figure.
\end{proof}
\begin{definition}[Matches Handle Geometry]
  \label{def:physics:phyx-mini-0823:phyxminiproblems-problemphyxmini0823-matcheshandlegeometry}
  \lean{PhyXMiniProblems.ProblemPhyXMini0823.MatchesHandleGeometry}
  \uses{def:physics:phyx-mini-0823:phyxminiproblems-problemphyxmini0823-lengthinmeters, def:physics:phyx-mini-0823:phyxminiproblems-problemphyxmini0823-positioninmeters, def:physics:phyx-mini-0823:phyxminiproblems-problemphyxmini0823-xhat, def:physics:phyx-mini-0823:phyxminiproblems-problemphyxmini0823-yhat, def:physics:phyx-mini-0823:phyxminiproblems-problemphyxmini0823-pipefittingtorquesetup}
  Cartesian realization of the pictured handle. From the fitting center the force point lies up-left: its horizontal component is negative and its vertical component is positive. This is geometry only, not a torque result
\end{definition}
\begin{proof}
  This is the declared model definition of Matches Handle Geometry.
\end{proof}
\begin{definition}[Matches Downward Force Geometry]
  \label{def:physics:phyx-mini-0823:phyxminiproblems-problemphyxmini0823-matchesdownwardforcegeometry}
  \lean{PhyXMiniProblems.ProblemPhyXMini0823.MatchesDownwardForceGeometry}
  \uses{def:physics:phyx-mini-0823:phyxminiproblems-problemphyxmini0823-forceinnewtons, def:physics:phyx-mini-0823:phyxminiproblems-problemphyxmini0823-forcemagnitudeinnewtons, def:physics:phyx-mini-0823:phyxminiproblems-problemphyxmini0823-yhat, def:physics:phyx-mini-0823:phyxminiproblems-problemphyxmini0823-pipefittingtorquesetup}
  Cartesian force direction corresponding to the vertical red arrow
\end{definition}
\begin{proof}
  This is the declared model definition of Matches Downward Force Geometry.
\end{proof}
\begin{definition}[Satisfies Torque Law]
  \label{def:physics:phyx-mini-0823:phyxminiproblems-problemphyxmini0823-satisfiestorquelaw}
  \lean{PhyXMiniProblems.ProblemPhyXMini0823.SatisfiesTorqueLaw}
  \uses{def:physics:phyx-mini-0823:phyxminiproblems-problemphyxmini0823-vectorreadout, def:physics:phyx-mini-0823:phyxminiproblems-problemphyxmini0823-spatialcross, def:physics:phyx-mini-0823:phyxminiproblems-problemphyxmini0823-pipefittingtorquesetup}
  The governing moment-of-force law in every coherent unit system. The lever arm is measured from the fitting center to the force-application point. This law contains no numerical torque answer
\end{definition}
\begin{proof}
  This is the declared model definition of Satisfies Torque Law.
\end{proof}
\begin{definition}[Answer Choice]
  \label{def:physics:phyx-mini-0823:phyxminiproblems-problemphyxmini0823-answerchoice}
  \lean{PhyXMiniProblems.ProblemPhyXMini0823.AnswerChoice}
  Labels of the four multiple-choice magnitudes
\end{definition}
\begin{proof}
  This is the declared model definition of Answer Choice.
\end{proof}
\begin{definition}[displayed Torque Magnitude In Newton Meters]
  \label{def:physics:phyx-mini-0823:phyxminiproblems-problemphyxmini0823-displayedtorquemagnitudeinnewtonmeters}
  \lean{PhyXMiniProblems.ProblemPhyXMini0823.displayedTorqueMagnitudeInNewtonMeters}
  \uses{def:physics:phyx-mini-0823:phyxminiproblems-problemphyxmini0823-answerchoice}
  Torque magnitude printed beside each answer label, in newton-metres
\end{definition}
\begin{proof}
  This is the declared model definition of displayed Torque Magnitude In Newton Meters.
\end{proof}
\begin{definition}[recorded Answer Choice]
  \label{def:physics:phyx-mini-0823:phyxminiproblems-problemphyxmini0823-recordedanswerchoice}
  \lean{PhyXMiniProblems.ProblemPhyXMini0823.recordedAnswerChoice}
  \uses{def:physics:phyx-mini-0823:phyxminiproblems-problemphyxmini0823-answerchoice}
  The answer label recorded in the dataset metadata
\end{definition}
\begin{proof}
  This is the declared model definition of recorded Answer Choice.
\end{proof}
\begin{definition}[Agrees At Ten Newton Meter Precision]
  \label{def:physics:phyx-mini-0823:phyxminiproblems-problemphyxmini0823-agreesattennewtonmeterprecision}
  \lean{PhyXMiniProblems.ProblemPhyXMini0823.AgreesAtTenNewtonMeterPrecision}
  \uses{def:physics:phyx-mini-0823:phyxminiproblems-problemphyxmini0823-torquemagnitudeinnewtonmeters, def:physics:phyx-mini-0823:phyxminiproblems-problemphyxmini0823-pipefittingtorquesetup, def:physics:phyx-mini-0823:phyxminiproblems-problemphyxmini0823-answerchoice, def:physics:phyx-mini-0823:phyxminiproblems-problemphyxmini0823-displayedtorquemagnitudeinnewtonmeters}
  Agreement with a whole-ten newton-metre display value: the exact magnitude is within half of one 10 N m display increment
\end{definition}
\begin{proof}
  This is the declared model definition of Agrees At Ten Newton Meter Precision.
\end{proof}
\begin{definition}[torque Sense]
  \label{def:physics:phyx-mini-0823:phyxminiproblems-problemphyxmini0823-torquesense-value}
  \lean{PhyXMiniProblems.ProblemPhyXMini0823.torqueSense}
  \uses{def:physics:phyx-mini-0823:phyxminiproblems-problemphyxmini0823-torquevectorquantity, def:physics:phyx-mini-0823:phyxminiproblems-problemphyxmini0823-torqueinnewtonmeters, def:physics:phyx-mini-0823:phyxminiproblems-problemphyxmini0823-torquesense}
  Rotation sense read from the sign of the out-of-page torque component
\end{definition}
\begin{proof}
  This is the declared model definition of torque Sense.
\end{proof}
% --- Archon declaration topology end ---
% --- Archon physics formalization source end ---
